\documentclass[answers]{exam}
\usepackage{../../mypackages}
\usepackage{../../macros}

\title{Savoir chercher l'information par vous-mêmes}
\author{N. Bancel}
\date{28 Novembre 2025}

% Mise en forme des solutions en bleu \SolutionEmphasis{\color{blue}} \renewcommand{\solutiontitle}{\noindent}

\begin{document}

\textbf{Collège Lycée Suger}
\hfill
\textbf{Mathématiques} \\

\textbf{Année 2025-2026}
\hfill
\textbf{Classe de 6ème} \par

{\let\newpage\relax\maketitle}

\begin{center} \textbf{\textcolor{red}{Ce document vous explique comment devenir autonomes pour chercher une information ou comprendre une notion.}} \ \end{center}

\section{Chercher dans le manuel}

Lorsque vous ne comprenez pas une notion, un mot scientifique ou une méthode :

\begin{compactitem}
  \item Commencez par regarder dans le manuel de mathématiques.
  \item Le manuel « lelivrescolaire » est disponible sur \textbf{ÉcoleDirecte}. Vous pouvez y retrouver :
  \begin{compactitem}
    \item les définitions importantes ;
    \item des exemples ;
    \item des méthodes expliquées pas à pas.
    \item des exercices corrigés portant sur la notion que vous ne comprenez pas.
  \end{compactitem}
\item Utilisez l’index ou les titres des chapitres pour retrouver rapidement la notion recherchée.
\end{compactitem}

\section{Chercher dans l’espace de travail}

Avant de poser une question, vérifiez si la réponse n’a pas déjà été donnée. Si elle porte sur 

\begin{compactitem}
  \item \textbf{Une interrogation / devoir sur table} : Regardez dans l’espace de Travail, dans le dossier \textbf{Mathématiques - 6èmes > Interrogations} et consultez les corrigés détaillés 
  \item \textbf{Une notion mal comprise d'un chapitre / une définition mal notée}. Par exemple, la définition ou la compréhension de la notion de médiatrice dans le chapitre 7. Allez dans Regardez dans l’espace de Travail, 
    dans le dossier \textbf{Mathématiques - 6èmes > [Chapitre 7] Droites, segments, demi-droites} et consultez
    (1) soit le cours disponible en ligne (2) soit le document des "vidéos pour mieux intégrer le cours". Ce sont des liens vers des vidéos explicatives.
  \item \textbf{Une méthode de calcul mental ou un corrigé de calcul mental} : les méthodes sont dans \textbf{Mathématiques - 6èmes > [Chapitre 0] Calcul mental} et les interrogations et leurs corrigés sont dans \textbf{Mathématiques - 6èmes > [Chapitre 0] Calcul mental > Interrogations}
\end{compactitem}

\section{Chercher une réponse sur Internet}

Vous pouvez utiliser Internet, mais il faut le faire intelligemment. Voici une méthode :

\subsection{Chercher avec des mots-clés simples}

\begin{compactitem}
  \item Utilisez Google ou Youtube.
  \item Précisez votre niveau : ajoutez "6ème" dans la recherche.
  \item Tapez une phrase courte : "exemple exercices corrigés sur médiatrice 6ème", "fraction 6ème cours", "périmètre rectangle explication 6ème", "reconnaître angle droit 6ème", "tracer médiatrice 6ème".
\end{compactitem}

\subsection{Lire la documentation (quand on code : GeoGebra, Python, Scratch)}

Personne n'est "bilingue" en code, et sait répondre à n'importe quelle question relative à une façon de tracer, programmer quelque chose. Lorsque vous utilisez un logiciel ou un langage de programmation, il est souvent très utile de se référer la documentation officielle. 

\vspace{1em}

\textbf{Exemple de situation} \\

Sur Geogebra, vous voulez tracer un angle de 33° dont les deux cotés font la même longueur (4 cms). Vous ne savez pas comment faire. Ici, la mesure de l'angle ou la longueur n'a pas d'importance - ce qui est important c'est de comprendre la méthode, quelles que soient les valeurs des angles et longueur.
Et ce qui importe aussi, c'est l'outil que vous utilisez : Geogebra. \\

Il est souvent plus efficace de faire la recherche en Anglais mais vous pouvez aussi le faire en Français. \\

\begin{compactitem}
  \item Cherchez "[Nom de l'outil] how to [ce que vous voulez faire]" ou "[Nom de l'outil] comment [ce que vous voulez faire]".
  \item Donc ici : "Geogebra how to draw angle with given length sides" ou "Geogebra comment tracer un angle avec des côtés de longueur donnée".
\end{compactitem}

\vspace{1em}

3 solutions

\begin{compactitem}
  \item Vous cliquez sur le premier lien (qui est souvent le lien considéré par Google comme le plus pertinent (qui correspond précisément à votre question))
  \item Vous cliquez sur le lien de la documentation officielle de Geogebra : \href{Geometry / Angles}{https://www.geogebra.org/math/angles?lang=en}
  \item Vous cliquez des vidéos YouTube recommandées qui semblent correspondre à votre recherche
\end{compactitem} 

\vspace{1em}

Servez-vous de cette méthode pour résoudre ces deux exercices entièrement par vous-mêmes, simplement en cherchant l'information sur Internet. 

\begin{figure}[H]
  \centering
  \includegraphics[width=0.5\textwidth]{figure_1.jpg}
  \caption{Exercice Geogebra}
\end{figure}

\begin{figure}[H]
  \centering
  \includegraphics[width=0.5\textwidth]{figure_2.jpeg}
  \caption{Exercice Geogebra}
\end{figure}


\subsection{Utiliser des sites fiables}

Voici quelques sites fiables pour les mathématiques :
\begin{compactitem}
  \item \textbf{Sesamath} : cours, exercices interactifs. Voici le lien vers le Manuel de 6ème : \url{https://manuel.sesamath.net/numerique/?ouvrage=cm6_2025}
  \item \textbf{Les vidéos d’Yvan Monka} (YouTube) : explications claires, adaptées aux collégiens. 
    \begin{compactitem}
    \item Chaîne YouTube : \url{https://www.youtube.com/channel/UCaDqmzanCq4ZYhdEm0Df9Qg}
    \item Le lien vers ses cours écrits de 6ème : \url{https://www.maths-et-tiques.fr/index.php/cours-maths/cours-de-maths-niveau-sixieme}
    \end{compactitem}
  \item \textbf{Mathenpoche} : exercices progressifs. Niveau 6ème : \url{https://mathenpoche.sesamath.net/?page=sixieme}
  \item Sites institutionnels : Lumni niveau 6ème \url{https://www.lumni.fr/college/sixieme}
\end{compactitem}

\section{Utiliser une IA}

Une IA peut vous aider à comprendre une notion ou à vous guider dans un exercice, mais il faut savoir bien formuler votre demande.

\subsection*{Comment poser une bonne question ?}

Une bonne question contient :
\begin{compactitem}
  \item la notion que vous souhaitez comprendre ;
  \item votre niveau (6ème) ;
  \item ce que vous avez déjà essayé ;
  \item l’endroit précis où vous bloquez.
\end{compactitem}

\subsection*{Exemples de bonnes questions}

\begin{compactitem}
  \item "Je suis en 6ème. Je veux comprendre comment calculer le périmètre d’un rectangle. Je ne comprends pas pourquoi on additionne deux fois la longueur et deux fois la largeur. Peux-tu te mettre dans la peau d'un professeur de 6ème et m'expliquer pas à pas ?"
  \item "J'ai un exercice sur les fractions. Je comprends comment les lire mais je ne sais pas comment comparer deux fractions. Peux-tu te mettre dans la peau d'un professeur de 6ème et m'expliquer la méthode ?"
\end{compactitem}

\section{Comment poser une question au professeur ?}

Pour que je puisse vous aider efficacement, votre question doit être \textbf{claire et structurée}. Voici le modèle :

\begin{compactitem}
  \item \textbf{Expliquer votre problème} : Voilà ce que je n'ai pas compris et qui me pose problème.
  \item \textbf{Dire ce que vous avez déjà essayé} : Voilà comment j'ai cherché (manuel, espace de travail, Internet)"
  \item \textbf{Essayer de décrire précisément le blocage}. Ce n'est pas toujours possible, parfois vous n'arriverez pas à formuler ce qui vous bloque. Mais essayez quand même.
  \item \textbf{Etre précis sur comment je peux (moi le professeur) retrouver l'information}. Si vous mentionnez un exercice, je voudrais le N° de l'exercice et la page. 
\end{compactitem}

\vspace{1em}

Ce type de question me permet de comprendre exactement où vous en êtes, et je peux vous aider beaucoup plus rapidement. \\

Par ailleurs, essayez de mettre un peu les formes : un "bonjour" au début, un "s'il vous plaît", un "merci" à la fin, et aussi une réponse pour me tenir
au courant une fois que j'ai répondu à votre email pour savoir si ça vous a vraiment aidé ou si vous êtes toujours bloqué, ça fait toujours plaisir !

\end{document}
